\newpage
\section{Задание 3}

\subsection*{Постановка задачи}

Расчет времени передачи файлов.
Сколько времени потребуется на передачу файла размером 83 Мбайт,
при размере кадра канального уровня 988 байт и доле повторно
передаваемых пакетов $7{,}81/1000$.

\begin{itemize}
    \item В качестве протокола канального уровня используется
    Ethernet DIX, на сетевом уровне --- IPv4, размеры заголовков
    транспортного и прикладного уровней вместе составляют 26 байт.
    Физический уровень --- 100BaseT.

    \item Время на передачу одного пакета данных составляет 89 мс.
    Задержка при повторной передаче пакета составляет 300 мс без
    учета времени повторной передачи.

    \item Повторная передача осуществляется после передачи
    $1000/7{,}81$ пакетов с округлением в большую сторону до
    ближайшего целого.
\end{itemize}

\begin{table}[h!]
\centering
\begin{tabular}{|c|l|}
\hline
\textbf{Параметр} & \textbf{Значение} \\ \hline
A & 83 Мбайт \\ \hline
B & 988 Байт \\ \hline
C & 7,81 \\ \hline
\end{tabular}
\caption{Исходные данные параметров A, B и C (вариант 22)}
\label{tab:abc_values}
\end{table}

\subsection*{Ход выполнения}

Размер кадра канального уровня составляет 988 байт. В состав кадра
входят заголовок Ethernet DIX размером 14 байт, заголовок IPv4
размером 20 байт, заголовки транспортного и прикладного уровней
общим размером 26 байт, а также данные файла.

\begin{table}[H]
\centering
\begin{tabular}{|c|c|c|c|}
\hline
\multicolumn{4}{|c|}{\textbf{Размер кадра --- 988 байт}} \\ \hline
\textbf{Ethernet DIX} &
\textbf{IPv4} &
\textbf{TP + APP} &
\textbf{Данные файла} \\ \hline
14 байт &
20 байт &
26 байт &
\textbf{928 байт} \\ \hline
\end{tabular}
\caption{Распределение заголовков и полезной нагрузки кадра}
\label{tab:packet_structure}
\end{table}

Полезная нагрузка одного кадра определяется как разность между
размером кадра и суммарным размером заголовков:

\[
D=988-14-20-26=928\text{ байт}.
\]

Таким образом, в одном кадре передается:

\[
\boxed{928\text{ байт данных файла}}
\]

Размер файла:

\[
F=83\text{ Мбайт}
=83\cdot1024^2
=87\,031\,808\text{ байт}.
\]

Количество кадров, необходимое для передачи файла, определяется
делением размера файла на размер полезной нагрузки одного кадра.
Поскольку количество кадров должно быть целым, результат
округляется в большую сторону:

\[
N=
\left\lceil
\frac{87\,031\,808}{928}
\right\rceil
=
93\,785.
\]

Таким образом, для передачи файла потребуется:

\[
\boxed{93\,785\text{ кадров}}
\]

Время передачи одного пакета составляет:

\[
t_p=89\text{ мс}.
\]

Доля повторно передаваемых пакетов составляет:

\[
\frac{7{,}81}{1000}.
\]

Согласно условию, повторная передача осуществляется после каждых

\[
\frac{1000}{7{,}81}\approx128{,}04
\]

пакетов. С учетом требования округления в большую сторону получаем:

\[
P=
\left\lceil
\frac{1000}{7{,}81}
\right\rceil
=129.
\]

Следовательно, повторная передача осуществляется после каждых
129 переданных пакетов.

Задержка при повторной передаче составляет 300 мс и не включает
время самой повторной передачи. Поэтому одна повторная передача
добавляет:

\[
t_{\text{повт}}=300+89=389\text{ мс}.
\]

Время обычной передачи всех кадров составляет:

\[
T_{\text{осн}}
=
93\,785\cdot89
=
8\,346\,865\text{ мс}.
\]

В секундах:

\[
T_{\text{осн}}
=
8\,346{,}865\text{ с}.
\]

Количество повторных передач определяется следующим образом:

\[
R=
\left\lfloor
\frac{93\,785-1}{129}
\right\rfloor
=
727.
\]

Следовательно, будет выполнено:

\[
\boxed{727\text{ повторных передач}}
\]

Дополнительное время, затрачиваемое на повторные передачи:

\[
T_{\text{повт}}
=
727\cdot389
=
282\,803\text{ мс}.
\]

В секундах:

\[
T_{\text{повт}}
=
282{,}803\text{ с}.
\]

Полное время передачи файла:

\[
T=
T_{\text{осн}}+T_{\text{повт}},
\]

\[
T=
8\,346\,865+282\,803
=
8\,629\,668\text{ мс}.
\]

Переведем результат в секунды:

\[
T=
8\,629{,}668\text{ с}.
\]

В минутах:

\[
T=
\frac{8\,629{,}668}{60}
\approx143{,}828\text{ мин}.
\]

Переведем результат в часы, минуты и секунды:

\[
\boxed{
T\approx2\text{ ч }23\text{ мин }50\text{ с}
}
\]

Таким образом, для передачи файла размером 83 Мбайт потребуется
примерно:

\[
\boxed{8\,629{,}668\text{ с}\approx2\text{ ч }23\text{ мин }50\text{ с}}
\]